Valuing changes in purchases and predicting where people choose to live require
different restrictions. Demand analysis begins with the expenditure needed to
attain a utility level. Its price derivatives connect purchases to the cost of
that utility. The Almost Ideal Demand System turns this relationship into a
specified system of budget shares, with restrictions governing adding up,
homogeneity, and symmetry \cite{deaton_muellbauer_1980}. This provides a useful
way to connect preference assumptions to an estimated response, rather than
treating every spending change as a change in income.

The distributional question adds differences in what consumers buy. When initial
baskets differ, the same relative price change can have different welfare effects.
The unequal-gains approach develops this expenditure channel within a specified
demand and trade structure \cite{fajgelbaum_khandelwal_2016}. Its restrictions
permit a mapping from aggregate shares to incidence across expenditure levels;
the mapping is not a general consequence of observing those shares. Earnings
incidence remains a separate part of the welfare account.

Endogenous location choices raise another question: how does policy change the
allocation that generates prices, earnings, and amenities? Spatial-policy analysis
defines a competitive allocation and the corresponding feasible planning problem,
then studies transfers in an economy with spillovers and sorting
\cite{fajgelbaum_gaubert_2020}. Moving from initial-basket valuation to this
policy comparison therefore requires an equilibrium account of adjustment and
ownership. These approaches share an interest in welfare, but the additional
structure answers an additional economic question; it does not independently
validate income recovery from a bilateral spending regression.
